\section{Conclusion}
\label{sec:conclusion}

This study evaluates the causal effect of China's 2017
climate-adaptive city pilot programme on urban ecological resilience
using a Double Machine Learning framework with satellite-derived
outcomes. Drawing on MODIS NDVI, ERA5 extreme weather events, and
municipal yearbook data for 338 cities over 2005--2023, we construct
a Climate-Stress Ecological Elasticity index that directly measures
ecosystem resistance to and recovery from climate shocks. The
principal findings are threefold.

First, the climate-adaptive city pilot has a positive, statistically
significant effect on ecological resilience, with the XGBoost DML
estimate of $\theta = 0.040$ ($p < 0.05$) representing a 5.3\%
increase in CSEE relative to the sample mean. This effect is robust
to alternative ML learners, 200-iteration placebo tests, propensity-score
matching, and subsample restrictions, with the post-2010 and post-2008
subsamples both achieving 5\% significance under Random Forest.

Second, the policy's effect operates almost entirely through the
recovery channel: the recovery component (RC) shows a highly
significant effect ($\theta = 0.081$, $p < 0.01$), while the
resistance component (CR) is insignificant. This decomposition
reveals that the pilot programme enhances resilience not by buffering
vegetation during extreme events but by accelerating post-disturbance
recovery---a mechanism consistent with the policy's emphasis on
ecological restoration infrastructure.

Third, the parallel-trends assumption is satisfied for all five
outcome variables, and the heterogeneity analysis reveals a
consistent positive direction across seven dimensions, with eastern
and smaller cities benefiting most. The comparison with an
alternative LST outcome, which fails the parallel-trends test,
demonstrates the importance of policy-specific identification and
parallel-trends validation.

These findings carry implications for both environmental management
practice and causal inference methodology. For practitioners, the
recovery-dominated mechanism suggests that climate-adaptation
investments should prioritise post-disturbance restoration
infrastructure over real-time protective measures, and that programme
scaling should target regions with strong implementation capacity.
For methodologists, the study demonstrates that DML--DID can deliver
meaningful efficiency gains over conventional TWFE-DID---approximately
25\% smaller standard errors---while maintaining the causal
interpretation of the estimate, and that algorithm comparison across
multiple ML learners serves as a valuable specification robustness
check in environmental policy evaluation.
